\chapter{Fundamentação Teórica}
\label{cap:fundamentacao}

Este capítulo reúne os fundamentos sobre os quais a arquitetura do Capítulo \ref{cap:arquitetura} se apoia, dos paradigmas de modelagem urbana aos padrões de engenharia de software para processamento assíncrono e persistência especializada.

\section{Cidades Inteligentes e Sociedade 5.0}
\label{sec:cidades-inteligentes}

Cidades inteligentes usam tecnologias digitais para otimizar recursos, melhorar os serviços públicos e elevar a qualidade de vida da população. Esses sistemas quase nunca são aplicações isoladas: eles integram subsistemas que pertencem a domínios administrativos diferentes. \citeonline{cavalcante2017challenges} destacam, entre as principais dificuldades dessa composição, a heterogeneidade tecnológica e a necessidade de permitir a evolução independente dos subsistemas. Em iniciativa municipal brasileira, integrar fontes de diferentes secretarias foi o principal obstáculo técnico \cite{cacho2016smart}. Trabalhos anteriores atacaram esse obstáculo por caminhos convergentes: um \textit{middleware} organizado em camadas geográficas \cite{souza2018geographic}, consolidado depois em plataforma para dados heterogêneos e aplicações urbanas \cite{pereira2022platform}, a publicação de dados abertos municipais \cite{macedo2017comparative} e a integração de sistemas de centros de operações urbanas \cite{loss2024integrating}.

A Sociedade 5.0 é um conceito originado na política japonesa de ciência e tecnologia, discutido por \citeonline{deguchi2020from}, que propõe colocar o bem-estar humano no centro do desenvolvimento tecnológico. É uma mudança em relação à Indústria 4.0, cujo foco está na eficiência dos processos produtivos. Para a engenharia de software urbano, pensar o sistema a partir do cidadão significa que ele deixa de ser fonte de dados agregados e passa a ser uma entidade central do modelo, com identidade, atributos e relacionamentos próprios. A crítica a essa mudança atinge dois pontos. O primeiro é que adotar o discurso centrado no humano não garante participação efetiva \cite{taylor2023human,tappert2024citizen}. O segundo vem de \citeonline{dejaeger2026right}, que demonstram que as decisões tomadas na camada arquitetural condicionam as formas de participação possíveis depois, e muitas vezes as restringem. É um argumento para tratar o \textit{backend} como objeto de pesquisa, e não como detalhe de implementação.

\section{Gêmeos Digitais: do Conceito Industrial à Cidade}
\label{sec:gemeos-digitais}

Gêmeos digitais são representações virtuais de objetos, sistemas ou processos reais, usadas para monitoramento, análise e simulação contínuos. \citeonline{grieves2016digital} formalizaram o conceito como o conjunto de construções de informação virtual que descreve um objeto físico ao longo de todo o seu ciclo de vida. Essa definição destaca dois aspectos que seguem centrais: a descrição completa do objeto e o acompanhamento de suas mudanças ao longo do tempo, e não em um instante isolado.

Com a evolução do conceito surgiram interpretações diferentes sobre o nível de fidelidade da representação, a sincronização com o objeto físico e a possibilidade de atuar sobre ele \cite{fuller2020digital}. \citeonline{vanderhorn2021digital} descrevem o gêmeo digital a partir de três elementos: o sistema físico, seu modelo virtual e a conexão entre ambos. Quando a conexão conta como componente, e não como subproduto, três níveis se distinguem: o \textit{modelo digital} não recebe dados automaticamente do objeto físico, a \textit{sombra digital} recebe dados em apenas uma direção e o gêmeo digital troca informação nos dois sentidos. Esta pesquisa parte desse último nível e avança um passo: além de fechar o ciclo de atuação, busca fazer esse ciclo raciocinar sobre o contexto, de modo que a resposta emitida dependa de quem a recebe, e não só do evento detectado.

Aplicar gêmeos digitais à cidade traz desafios que vão além da complexidade técnica. Diferentemente de um equipamento industrial, a cidade não tem um modelo único de referência contra o qual o gêmeo possa ser verificado, e seus processos constitutivos são sociais \cite{batty2018digital}. \citeonline{yeh2024from} reconstroem a evolução da modelagem urbana até o gêmeo digital.

Como registrado na Seção \ref{sec:contextualizacao}, muitas iniciativas municipais ainda não alcançam o conceito completo. Grande parte funciona como sombra digital, recebendo informação do ambiente sem fechar o ciclo de atuação \cite{ferrebigorra2022adoption}. Os níveis de maturidade em que o sistema apoia a decisão e executa intervenções seguem raros, e isso decorre de limitações arquiteturais e institucionais, não da falta de sensores \cite{haraguchi2024assessing}. Estudos recentes já consideram a dimensão humana nas etapas de percepção e análise \cite{luo2025perception,qanazi2025integrating}, mas permanece a lacuna na representação dos perfis individuais e no uso dessa informação pelo \textit{backend} durante a execução dos serviços. Esta pesquisa se concentra nessa lacuna.

\section{Gêmeos Digitais Humanos e Hiperpersonalização}
\label{sec:hdt}

Considerar a infraestrutura urbana e o cidadão no mesmo sistema permite sair das respostas genéricas. No contexto industrial essa abordagem já está sistematizada, incluindo soluções que usam computação de borda para manter a comunicação entre a pessoa e seu gêmeo virtual \cite{asad2023human,okegbile2022edge}. Nesses cenários, a latência do canal entra como requisito de projeto.

\citeonline{fontes2024human} examinam em que medida gêmeos digitais humanos poderiam habilitar os princípios da Sociedade 5.0 e identificam três riscos: a assimetria de informação entre quem opera o sistema e quem é representado por ele, a discriminação algorítmica derivada de perfis inferidos e a dificuldade de assegurar consentimento informado sobre inferências que o próprio indivíduo não antecipa. Deles decorrem os requisitos de minimização de dados, explicabilidade e revogabilidade do consentimento adotados no Capítulo \ref{cap:arquitetura}.

Hiperpersonalização é adaptar serviços às características de cada indivíduo em granularidade fina e em tempo real, em vez de usar categorias gerais de usuário \cite{desai2022hyper}. \citeonline{lex2025hybrid} apresentam abordagens que combinam modelos estatísticos e conhecimento simbólico, direção próxima à adotada aqui, em que a personalização resulta de regras aplicadas a um grafo semântico, e não de um modelo opaco. No domínio urbano, \citeonline{liu2023urban} demonstram o uso de grafos de conhecimento urbanos para construir perfis de usuários móveis, e \citeonline{wins2024enhancing} propõem a personalização de itinerários em múltiplos estágios.

A abordagem também é relevante em situações de emergência. \citeonline{hashemi2018emergency} revisa simulações, exercícios e modelos de evacuação de pessoas com deficiência e documenta a lacuna de desempenho desse grupo. Outros trabalhos mostram que considerar restrições individuais de acessibilidade é tecnicamente viável em roteamento e em mobilidade urbana \cite{deng2022semantic,allahbakhshi2026zuriact}.

\section{Sensibilidade ao Contexto e Interoperabilidade Semântica}
\label{sec:contexto}

\citeonline{dey2001understanding} define contexto como qualquer informação utilizável para caracterizar a situação de uma entidade, e considera um sistema sensível ao contexto quando ele usa essa informação para oferecer informação ou serviço relevante ao usuário. Daí decorrem duas implicações de projeto: a relevância depende da tarefa de quem recebe a resposta, e a sensibilidade ao contexto é propriedade do comportamento do sistema, não de sua base de dados.

Os estudos sobre \textit{middleware} sensível ao contexto abordam principalmente a coleta, a representação, o processamento e a distribuição dessas informações \cite{liang2015social,sezer2018context,lopes2009context,lopes2013middleware}. Para esta pesquisa, três aspectos são importantes: manter um modelo de contexto compartilhado entre os componentes, permitir que mudanças sejam propagadas com baixa latência, o que exige que os componentes não fiquem esperando uns pelos outros; e separar a representação dos dados das regras utilizadas para interpretá-los. Essa separação facilita a evolução das regras sem exigir alterações constantes no modelo de dados.

Em uma cidade inteligente, informação sobre clima, segurança pública, mobilidade e perfis de usuário precisa ser cruzada para que o sistema reconheça uma situação. Para isso os dados têm de manter significado consistente mesmo vindo de fontes distintas. A arquitetura ancora-se, por essa razão, no padrão NGSI-LD, especificado pelo \textit{European Telecommunications Standards Institute} \cite{etsi2024ngsild}.

O padrão anterior, o NGSI-v2, usava estruturas planas em \textit{JavaScript Object Notation} (JSON) e nomes de atributo em texto livre, o que gerava ambiguidade na integração de bases diferentes. Um atributo chamado \texttt{temperatura} tanto pode ser a medição do ar na rua quanto o estado térmico de um equipamento, e nada no dado permite distinguir os dois casos. Os relacionamentos entre entidades também dependiam de convenção de nome, com prefixos como \texttt{ref}, que um programa lê mas não interpreta como ligação.

O NGSI-LD resolve esses dois pontos ao incorporar os princípios da Web Semântica e representar o domínio como um grafo de entidades e relacionamentos. Quatro recursos importam aqui: a identificação global por \textit{Uniform Resource Identifier} (URI), o atributo \texttt{@context}, o relacionamento como dado de primeira classe e a associação de propriedades a outras propriedades. A identificação global das entidades por URI previne colisão de identidade em bases federadas. O atributo \texttt{@context} associa os nomes locais de atributo a conceitos de vocabulários compartilhados e, com isso, resolve a ambiguidade do exemplo acima. O relacionamento entre duas entidades é tratado como dado de primeira classe, no sentido definido no Capítulo \ref{cap:introducao}: o sistema registra não apenas onde um cidadão está, mas como ele se relaciona com um evento de emergência próximo. A travessia de grafo empregada no Capítulo \ref{cap:arquitetura} depende desse recurso. Por fim, propriedades podem ser associadas a outras propriedades para registrar o momento em que um dado foi coletado e a sua origem. Essa possibilidade sustenta a representação de proveniência por atributo apresentada na Seção \ref{sec:modelagem}.

A relação do NGSI-LD com outras abordagens de interoperabilidade semântica é discutida na Seção \ref{sec:resultados-msl} \cite{viola2019mapping,gonzalez2024reshaping}. A iniciativa \textit{Smart Data Models} \cite{smartdatamodels2025} disponibiliza modelos padronizados por domínio, entre eles dispositivos, alertas, transporte e parâmetros ambientais, o que ajuda a manter a consistência dos dados entre implementações independentes e complementa os esforços mais amplos de interoperabilidade semântica na Internet das Coisas \cite{ganzha2017semantic,qamar2019smart}.

\section{O Ecossistema FIWARE e o Context Broker}
\label{sec:fiware}

O FIWARE é um ecossistema de componentes de código aberto voltado ao desenvolvimento de aplicações inteligentes. Seu elemento principal é o \textit{Context Broker}, responsável pelo ciclo de vida da informação de contexto \cite{cirillo2019standard}. O ecossistema também oferece componentes de conexão com dispositivos, de processamento e de gestão de identidade e autorização \cite{bauer2022fiware,salhofer2019building}. No contexto dos gêmeos digitais urbanos, \citeonline{bauer2021urban} estabelecem a correspondência entre entidades NGSI-LD e os elementos que compõem o gêmeo, referência direta para este trabalho, e a fundação mantenedora registra seu posicionamento sobre gêmeos digitais e espaços de dados em \cite{ditorrepadula2025smart}.

Entre as implementações de \textit{Context Broker} compatíveis com NGSI-LD, esta pesquisa adota o Stellio \cite{stellio2025docs}. Desenvolvido em Kotlin com o arcabouço Spring, ele foi projetado como um conjunto de microsserviços. A comparação com as demais implementações compatíveis, o Orion-LD e o Scorpio, e os motivos da escolha constam da Seção \ref{sec:adr}.

Sua arquitetura reúne serviços distintos. O \textit{API Gateway} recebe as requisições e encaminha cada chamada ao serviço correspondente. O serviço de busca atende às consultas de contexto e temporais sobre PostgreSQL com PostGIS e TimescaleDB. Já o serviço de assinaturas tem banco próprio e avalia as assinaturas que disparam notificação quando o contexto muda. A comunicação entre esses componentes passa por um barramento Kafka interno, no qual cada operação sobre o contexto gera um evento atômico. Esse registro de mudanças, somado à avaliação de assinaturas em serviço separado, permite acoplar componentes de domínio ao núcleo sem alterá-lo.

\section{Espaços de Dados, Soberania e Governança}
\label{sec:dataspaces}

A hiperpersonalização depende do acesso a atributos individuais, alguns deles sensíveis. Por isso, além de permitir o compartilhamento, a arquitetura precisa garantir que cada organização mantenha controle sobre o que disponibiliza. Essa é uma das propostas centrais dos \textit{Data Spaces}, ambientes em que organizações distintas compartilham dados de forma descentralizada, sob contratos de uso verificáveis \cite{otto2019international}.

O modelo de referência da \textit{International Data Spaces Association} \cite{idsa2022ram} estabelece os papéis e os processos desse arranjo e destaca o conector como ponto de aplicação das políticas de uso. O Gaia-X segue princípios semelhantes e leva o mesmo arranjo à escala europeia \cite{gaiax2022architecture,tardieu2022role}.

No contexto urbano, \citeonline{cuno2019data} discutem governança e soberania em espaços de dados municipais baseados em interfaces padronizadas. Outros estudos tratam de temas relacionados, como a soberania digital corporativa, a identidade auto-soberana e a convergência entre espaços de dados e gêmeos digitais \cite{sanchezserna2026data,somma2023digital,uslander2022symbiotic}. Também importa o custo desses mecanismos: \citeonline{galij2025evaluating} medem a sobrecarga de desempenho que a camada de soberania introduz, aspecto decisivo para sistemas que precisam responder rápido.

Além dos aspectos arquiteturais, o tratamento de dados pessoais responde a requisitos legais. Nesta pesquisa, esse processamento está sujeito à Lei Geral de Proteção de Dados Pessoais \cite{brasil2018lgpd}, com o regulamento europeu como referência complementar \cite{eu2016gdpr}, dada sua influência sobre os mecanismos de conformidade dos espaços de dados analisados. A gestão de identidade e autorização apoia-se em solução de código aberto para controle de acesso baseado em papéis \cite{keycloak2025}.

\section{Arquitetura Orientada a Eventos, Microsserviços e Persistência}
\label{sec:eda}

Em sistemas compostos por vários serviços, a comunicação síncrona cria dependência entre os componentes. Quando um serviço precisa aguardar a resposta de outro, qualquer lentidão afeta o funcionamento dos demais. O cenário piora em ambiente urbano, em que milhares de dispositivos publicam leituras ao mesmo tempo. Para lidar com isso, a arquitetura adota o modelo orientado a eventos.

Nesse modelo, toda mudança de estado relevante é registrada como evento imutável e disponibilizada aos componentes interessados. A detecção de fumaça por um sensor ou a alteração de rota de um cidadão gera um evento que será processado por quem o assinar. A comunicação segue o padrão \textit{Publish/Subscribe}, cujas variantes e propriedades foram sistematizadas por \citeonline{eugster2003many}: o produtor publica sem saber quem vai consumir, e o consumidor recebe apenas os eventos dos canais que acompanha. O ganho é o desacoplamento espacial, temporal e de sincronização. Assim, a indisponibilidade de um componente secundário, como um painel de visualização, não interrompe a ingestão de dados de emergência.

A arquitetura de microsserviços complementa essa abordagem ao dividir a aplicação em serviços de responsabilidade única, implantáveis e escaláveis de forma independente \cite{newman2021building}. O vocabulário de projeto adotado aqui vem do catálogo de \citeonline{richardson2018microservices}, em especial \textit{Backend for Frontend}, \textit{API Gateway} e \textit{Database per Service}. Na comunicação entre serviços, o Apache Kafka atua como plataforma de mensageria distribuída, concebida originalmente para processamento de \textit{logs} \cite{kreps2011kafka}. Sua abstração central é o \textit{log} particionado, ordenado e persistente, que dá durabilidade às mensagens, permite reproduzir o histórico de eventos e atende consumidores de velocidades distintas. \citeonline{kleppmann2017designing} relaciona esse modelo ao uso de \textit{logs} de eventos, visões materializadas e consistência eventual, conceitos adotados na organização da camada de persistência.

Um gêmeo digital urbano também lida com dois tipos de informação ao mesmo tempo. O estado corrente das entidades é consultado por proximidade, contenção e interseção espacial, enquanto o histórico de atributos é consultado por agregação em janelas temporais. Bancos unificados tendem a atender bem a um desses perfis em detrimento do outro, e esse desequilíbrio motiva a segregação.

Para a informação geográfica adota-se o PostgreSQL com a extensão PostGIS, que oferece tipos geométricos, índices espaciais e operações de análise \cite{obe2021postgis}. Para o histórico adota-se o TimescaleDB, que acrescenta ao PostgreSQL particionamento temporal, compressão e agregações contínuas \cite{timescale2025docs}. O uso de bancos de séries temporais em plataformas urbanas já foi tratado empiricamente \cite{john2023application,rudakov2023comparison}. Por fim, a conteinerização padroniza o ambiente de execução e permite reproduzir a implantação \cite{ahmed2018docker}.

A fundamentação reúne os conceitos e as tecnologias de que a proposta depende, mas não responde sozinha a uma questão: em que medida a literatura já reuniu essas peças em arquiteturas concretas de \textit{backend}. O capítulo seguinte investiga isso por um mapeamento sistemático com protocolo explícito.
